Please write a Mathematica program that computes the minimal number of contexts that carries a single 1 and three 0's (admissible two-valued state), and the remaining context with 4 vertices value 0's. Here is Cabello's 8-9 configuration (the first number is just D=4:




Pasting of 9 blocks
18 atoms
 9 blocks
 0 proper subsets of blocks
 4   1  2  3  4
 4   4  5  6  7
 4   7  8  9 10
 4  10 11 12 13
 4  13 14 15 16
 4  16 17 18  1
 4   2  9 11 18
 4   3  5 12 14
 4   6  8 15 17
0 2-valued evaluations of atoms:

set of 2-valued evaluations of atoms:
nonempty: no

############################################################################################


(*Cabello 8\[Dash]9 configuration Brute-force search (no Maximize/\
Reduce) to avoid recursion issues.Computes:maxAdmissible=maximal #\
 of contexts with exactly one 1 minZeroContexts=minimal # of all-\
zero contexts*)

ClearAll["Global`*"];

(*---1. Problem data---*)

atoms = Range[18];

blocks = {{1, 2, 3, 4}, {4, 5, 6, 7}, {7, 8, 9, 10}, {10, 11, 12,
    13}, {13, 14, 15, 16}, {16, 17, 18, 1}, {2, 9, 11, 18}, {3, 5, 12,
     14}, {6, 8, 15, 17}};

na = Length[atoms];
nb = Length[blocks];

(*---2. Brute-force search over all 2^18 valuations---*)

maxAdmissible = -1;
bestValuation =.;
bestBlockCounts =.;

(*Iterate over all {0,1}^18 assignments*)
Do[(*val is a list of 18 bits:val[[i]]=value of atom i*)
  val = assignment;
  (*For each block,count how many 1's appear in that block*)
  blockCounts = Total[val[[#]]] & /@ blocks;
  (*Valid if every block has either 0 or 1 ones*)
  If[AllTrue[blockCounts, (# == 0 || # == 1) &],
   nAdmissible = Count[blockCounts, 1];
   If[nAdmissible > maxAdmissible, maxAdmissible = nAdmissible;
    bestValuation = val;
    bestBlockCounts = blockCounts;];];, {assignment,
   Tuples[{0, 1}, na]}];

minZeroContexts = nb - maxAdmissible;

(*---3. Output---*)

Print["maxAdmissible = ", maxAdmissible];
Print["minZeroContexts = ", minZeroContexts];

Print["\nBest atom valuation v[i]:"];
Do[Print["  v[", i, "] = ", bestValuation[[i]]], {i, na}];

Print["\nContext counts (number of 1's in each block):"];
Do[Print["  Context ", j, "  ", blocks[[j]], "  ->  #1's = ",
   bestBlockCounts[[j]]], {j, nb}];

admissibleContexts = Flatten@Position[bestBlockCounts, 1];
zeroContexts = Flatten@Position[bestBlockCounts, 0];

Print["\nAdmissible contexts (exactly one 1): ", admissibleContexts];
Print["All-zero contexts: ", zeroContexts];



############################################################################################


(* 1. DEFINE CABELLO 18-9 CONFIGURATION *)
raw = {
   {4, 1, 2, 3, 4}, {4, 4, 5, 6, 7}, {4, 7, 8, 9, 10},
   {4, 10, 11, 12, 13}, {4, 13, 14, 15, 16}, {4, 16, 17, 18, 1},
   {4, 2, 9, 11, 18}, {4, 3, 5, 12, 14}, {4, 6, 8, 15, 17}
   };

(* Remove the dimension '4' from the input *)
blocks = Rest /@ raw;
atoms = Union[Flatten[blocks]];
vars = Table[x[i], {i, Length[atoms]}];

(* 2. DEFINE RULES *)
(* Rule: Single 1 (Sum = 1) *)
ruleOne[b_] := BooleanCountingFunction[{1}, Length[b]] @@ (x /@ b);
(* Rule: Four 0s (Sum = 0) *)
ruleZero[b_] := And @@ (! x[#] & /@ b);

(* 3. FIND MINIMAL SET FOR CONTRADICTION *)
Print["--- analyzing minimal contexts ---"];
minimalSet = {};

Do[
  subsets = Subsets[Range[Length[blocks]], {k}];
  (* Find a subset that CANNOT be colored with all 1s *)
  contradictionSubset = SelectFirst[subsets,
    SatisfiabilityInstances[And @@ (ruleOne[blocks[[#]]] & /@ #), vars] === {} &
  ];

  If[ListQ[contradictionSubset],
   minimalSet = contradictionSubset;
   Print["Minimal number of contexts required for contradiction: ", k];
   Print["Context indices: ", minimalSet];
   Break[];
   ];
  , {k, 1, Length[blocks]}
];

(* 4. COMPUTE AND OUTPUT VALUES FOR THE 'LEGAL' STATE *)
(* We seek a state where (k-1) blocks are '1', and exactly 1 block is '0' *)

If[Length[minimalSet] > 0,
  Print["\n--- Computing values for the 'Remaining Context 0' state ---"];

  (* We try to force the LAST context in the minimal set to be the '0' context,
     and the rest to be '1' contexts. *)

  targetZeroIndex = Last[minimalSet];
  standardIndices = Most[minimalSet];

  Print["Attempting to set Context #", targetZeroIndex, " to 0000, and others to 1000..."];

  equation = ruleZero[blocks[[targetZeroIndex]]] &&
             (And @@ (ruleOne[blocks[[#]]] & /@ standardIndices));

  solution = SatisfiabilityInstances[equation, vars];

  If[solution =!= {},
   (* Convert True/False to 1/0 *)
   values = Boole[First[solution]];

   Print["\nSolution found! Here are the vertex values as Mathematica code:\n"];

   (* PRETTY OUTPUT FOR COPY-PASTING *)
   Print["vertexValues = ", values, ";"];

   Print["\nInterpretation:"];
   Print["Atom 1 value: ", values[[1]]];
   Print["Atom 2 value: ", values[[2]]];
   Print["..."];

   (* VERIFICATION *)
   Print["\nVerification of Context sums:"];
   checkSum[ind_] := Total[values[[blocks[[ind]]]]];
   Do[
    val = checkSum[i];
    status = If[val == 1, " (OK: Single 1)", If[val == 0, " (OK: Four 0s)", " (INVALID)"]];
    Print["Context ", i, ": Sum = ", val, status];
    , {i, minimalSet}];
   ,
   Print["No such state exists for this specific choice of zero-context."]
   ];
 ];
